\documentclass[10pt,a4paper]{amsart}
\usepackage[margin=19mm]{geometry}
\usepackage{amsmath,amssymb,mathtools}
\usepackage[scr=rsfs,cal=euler]{mathalfa}
\usepackage{xfrac}
\usepackage[unicode,hidelinks,bookmarks=false]{hyperref}
\newcommand{\ie}{i.e.\ }
\newcommand{\cf}{cf.\ }
\newcommand{\resp}{respectively\ }
\newcommand{\Subsection}{\subsection}
\providecommand{\nobreakdash}{\nobreak\hbox{-}}
\setlength{\emergencystretch}{2em}
\multlinegap0pt
\usepackage{amssymb}
\usepackage{multirow}
\usepackage[all]{xy}
\usepackage{xcolor}
\newtheorem{lem}{Lemma}
\title{On the Euclidean Distance Degree of Quadratic Two-Neuron Neural Networks}
\author{Giacomo Graziani}
\date{Source: arXiv:2512.21016v3 (CC BY 4.0)}
\begin{document}
\maketitle

\noindent\emph{Source and licence.} On the Euclidean Distance Degree of Quadratic Two-Neuron Neural Networks. Version arXiv:2512.21016v3 (CC BY 4.0), distributed by arXiv under the Creative Commons Attribution 4.0 International licence, http://creativecommons.org/licenses/by/4.0/, as recorded in the arXiv licence metadata for this exact version and shown on the abstract page https://arxiv.org/abs/2512.21016v3. Full licence notice: \texttt{licenses/arxiv-2512.21016-CC-BY-4.0.txt}.\par\smallskip

\noindent\textit{Editorial scope.} Counted unit: the article's lem lem:BWInnerProductIndependent (source0.tex lines 361--366). The article's complete original proof of it is delivered with it (source0.tex lines 368--448, one \texttt{proof} environment of 81 lines). No mathematics is written, repaired or re-derived: the statement and the proof are the article's own lines, in the article's own notation, and no retained line is substituted or re-flowed.  No further source-local context is needed: every cross-reference of every retained line resolves inside the two delivered spans.  Nothing was omitted from any retained span, so the package carries no omission record.  No retained line contains a citation, so the wrapper carries no bibliography and no bibliography entry.  No retained line contains a figure environment, an image-inclusion command or drawing code, so no image is delivered and the package declares no asset dependency.  Editorial formatting only, in addition: an AMS \texttt{amsart} preamble that restates the article's own declarations and macros used by the retained lines, namely the environments \texttt{lem} on separate counters, together with the standard packages those lines need.  The retained environments print in this document as Lemma 1 in order of appearance, whereas the article prints the same items with the numbering of its own full document.  The complete native source of the article as distributed in the arXiv e-print for this exact version is bundled unchanged as \texttt{sources/191608/source0.tex}.
\begin{lem}
\label{lem:BWInnerProductIndependent}
The Bombieri-Weyl product is a scalar product, that is, it is
symmetric, linear in both variables, and positive-definite.
Moreover, it does not depend on the choice of the orthonormal basis.
\end{lem}

\begin{proof}
We include the proof for lack of reference. The first part is clear. To see that the Bombieri-Weyl scalar product
is invariant under orthogonal changes of basis let $Q$ be an
orthogonal transformation on $V$ and let
$f\in\mathrm{Sym}^{d}\left( V \right)$,
let moreover $n=\dim V$. Denote by $f\circ Q$ the polynomial defined
by
\[
\left( f\circ Q \right)\left( x \right)
=
f\left( Qx \right).
\]
Since the Veronese map is non-degenerate the set of powers of linear
forms $\left\{ v^{d}\mid v\in V \right\}$ spans the space
$\mathrm{Sym}^{d}\left( V \right)$. By the linearity of the scalar
product, it suffices to prove the invariance for elements of this
form. For $v\in V$ we identify it with the linear map it defines via
duality
\[
v\left( x \right)
=
\left\langle v,x \right\rangle_{V},
\]
hence the symmetric power is the polynomial
\[
v^{d}\left( x \right)
=
\left\langle v,x \right\rangle_{V}^{d}.
\]
Let $v,w\in V$ then applying the definition of the product and the
Multinomial Theorem:
\begin{equation}
\left\langle v^{d},w^{d} \right\rangle_{\mathrm{BW}}
=
\sum_{\left| \alpha \right|=d}
\binom{d}{\alpha}
v^{\alpha} w^{\alpha}
=
\left( \sum_{i=1}^{n} v_{i} w_{i} \right)^{d}
=
\left\langle v,w \right\rangle_{V}^{d}.
\label{eq:CompatibilityBWInnerProduct}
\end{equation}
Now consider the action of the orthogonal transformation $Q$. The
polynomial transforms as
\[
\left( v^{d} \right)\circ Q
=
\left( Q^{T}v \right)^{d}.
\]
Therefore:
\[
\left\langle
\left( v^{d} \right)\circ Q,
\left( w^{d} \right)\circ Q
\right\rangle_{\mathrm{BW}}
=
\left\langle
\left( Q^{T}v \right)^{d},
\left( Q^{T}w \right)^{d}
\right\rangle_{\mathrm{BW}}
=
\left\langle Q^{T}v,Q^{T}w \right\rangle_{V}^{d}.
\]
Since $Q$ is orthogonal,
\[
\left\langle Q^{T}v,Q^{T}w \right\rangle_{V}
=
\left\langle v,w \right\rangle_{V}.
\]
Thus
\[
\left\langle
\left( Q^{T}v \right)^{d},
\left( Q^{T}w \right)^{d}
\right\rangle_{\mathrm{BW}}
=
\left\langle v^{d},w^{d} \right\rangle_{\mathrm{BW}}.
\]
This proves the invariance.
\end{proof}

\end{document}
